\section{Photon Sphere Structure and Observable Deviation}

A central consequence of the TIG framework is the modification of null geodesic structure near compact objects. In particular, the location of the photon sphere provides a direct probe of the underlying gravitational geometry.

The photon sphere is determined by the standard condition for unstable circular null geodesics:
\begin{equation}
r F'(r) - 2F(r) = 0.
\end{equation}

Within the TIG framework, the metric function takes the form
\begin{equation}
F(r) = 1 - \frac{2M r^2}{r^3 + r_c^3},
\end{equation}
which introduces a characteristic scale $r_c$ governing deviations from classical General Relativity.

In the limit $r_c \to 0$, the model smoothly reduces to the Schwarzschild solution with
\begin{equation}
r_{\mathrm{ph}}^{\mathrm{GR}} = 3M,
\end{equation}
ensuring consistency with established gravitational theory.

We define the relative deviation from the Schwarzschild photon sphere as
\begin{equation}
\Delta(r_c) = \frac{r_{\mathrm{ph}}^{\mathrm{TIG}} - 3M}{3M}.
\end{equation}

Numerical solutions reveal a clear and systematic deviation from the classical result. While the deviation remains negligible for small $r_c$, it grows non-linearly as the characteristic scale increases, reaching values of approximately $-4%$ in the explored regime.

This shift implies a measurable contraction of the photon sphere and, consequently, a reduced black hole shadow radius relative to General Relativity.

Although current observational precision does not yet resolve deviations at this level, the predicted effect defines a concrete and testable target for future high-resolution observations.
